% 論文節の日本語訳（読み物 PDF）用 preamble
% lualatex + jlreq
\documentclass[12pt,paper=a4]{jlreq}
\usepackage{luatexja-fontspec}
\usepackage{amsmath,amssymb,bm}
\usepackage{graphicx}
\usepackage{hyperref}
\usepackage{enumitem}
\usepackage{booktabs}
\usepackage{microtype}
\usepackage{xcolor}
\usepackage{caption}
\usepackage{tcolorbox}
\tcbuselibrary{breakable,skins}

\hypersetup{
  colorlinks=true,
  linkcolor=blue!60!black,
  urlcolor=blue!60!black,
  citecolor=blue!60!black,
}

\graphicspath{{images/}{./}}

\setlength{\parindent}{1em}
\setlength{\parskip}{0.55em}
\setlist[itemize]{leftmargin=1.6em,itemsep=0.25em,topsep=0.35em}
\setlist[enumerate]{leftmargin=1.8em,itemsep=0.25em,topsep=0.35em}

\captionsetup{
  font=small,
  labelfont=bf,
  skip=0.4em,
  justification=raggedright,
  singlelinecheck=false,
}

\newcommand{\vx}{\bm{x}}
\newcommand{\vu}{\bm{u}}
\newcommand{\veps}{\bm{\epsilon}}
\newcommand{\vw}{\bm{w}}

\DeclareMathOperator*{\argmax}{arg\,max}
\DeclareMathOperator*{\argmin}{arg\,min}
\DeclareMathOperator{\Tr}{Tr}
\DeclareMathOperator{\diag}{diag}

% 命題・定理ボックス（ページ跨ぎ時も切れ目を抑える）
\newtcolorbox{propositionbox}[1][]{%
  enhanced,
  breakable,
  colback=blue!4!white,
  colframe=blue!45!black,
  coltitle=white,
  fonttitle=\bfseries\small,
  title={#1},
  attach boxed title to top left={yshift=-2mm, xshift=3mm},
  boxed title style={
    colback=blue!45!black,
    colframe=blue!45!black,
    boxrule=0pt,
    arc=1pt,
  },
  boxrule=0.5pt,
  arc=1.5pt,
  left=5pt, right=5pt, top=8pt, bottom=5pt,
  before skip=0.9em,
  after skip=0.9em,
  pad at break=3mm,
}

% 図はフロートせず本文位置に置く（分断を防ぐ）
\newcommand{\translationfigure}[2]{%
  \begin{center}
    \includegraphics[width=0.88\linewidth]{#1}\\[0.35em]
    {\small #2\par}
  \end{center}
  \vspace{0.4em}%
}

\begin{document}
{\color{gray}\footnotesize 出典: §8.4\par}
\vspace{0.2em}

\section*{8.4 訓練不要ガイダンス}


訓練不要ガイダンスは、拡散モデル自体を再訓練することなく、外部の条件・損失・測定目的を用いて固定された事前学習済み拡散モデルのサンプリングダイナミクスを修正する。式 (8.1.1) のBayes分解はスコア空間での有用な動機を与える：ガイダンスは無条件スコアに条件依存方向を加える。実務上、タスク固有損失は、予測クリーンサンプルまたはサンプリング中に用いるスコアへの補正を構成するために用いられ得る。
まず §8.4.1 でこれら2つの一般形を導入する。次いで §8.4.2 で、そのような代理誘導ダイナミクスが実際にどの分布を生成するかを問い、§8.4.3 で逆問題に対する代表的な訓練不要アプローチへ移る。

\textbf{設定と記法。} $c$ を条件付け変数とする。スコア予測で表された事前学習済み拡散モデル $s_{\phi^\times}(x_t,t)$ へのアクセスを仮定する\footnote{ここでは単純化のためスコアおよび $x$-予測パラメータ化を採用する；$\epsilon$-予測などの他のパラメータ化も同様に扱える。}。加えて、望ましい条件との不一致を測るタスク損失

\[
\ell(\cdot,c):\mathbb{R}^D\to\mathbb{R}
\]


が与えられているとし、$\ell(x,c)$ が小さいほど整合が良いことを示す。例には参照ターゲットへの知覚距離や、CLIP（Radford et al., 2021）のような事前学習表現から導かれる特徴空間距離または負の類似度が含まれる。
標準的な線形–ガウス前方ノイズ化核

\[
p_t(\cdot|x_0):=\mathcal{N}\bigl(\cdot;\alpha_t x_0,\sigma_t^2 I\bigr)
\]


を考える。具体例として、後に式 (9.2.2) で導入されるDDIM更新を借用する：

\[
x_{t\to t-1}=\alpha_{t-1}\underbrace{\hat{x}_0(x_t)}_{\text{データ空間}}-\sigma_{t-1}\sigma_t\underbrace{\hat{s}(x_t)}_{\text{ノイズ空間}}, \quad (8.4.1)
\]


ここで $\hat{x}_0(x_t):=x_{\phi^\times}(x_t,t)$ はクリーン $x$-予測、$\hat{s}(x_t):=s_{\phi^\times}(x_t,t)$ はノイズ水準 $t$ でのスコア予測である。


\subsection*{8.4.1 訓練不要ガイダンスの概念枠組み}


共通戦略は、タスク損失の勾配を用いて、サンプリング更新における予測クリーンサンプルまたはスコア項のいずれかを修正することである。これら2操作を区別するため、データ空間方向

\[
G_t^{\mathrm{data}}(x_t,c):=-\nabla_{x_0}\ell(x_0,c)\bigm|_{x_0=\hat{x}_0(x_t)}
\]


を定義する。ノイズ空間補正については、プロローグで導入したガイダンス方向

\[
G_t(x_t,c):=\nabla_{x_t}\log h_t(x_t;c)
\]


を再利用する。正のガイダンスポテンシャル $h_t$ は下記で指定する。
式 (8.4.1) のDDIM更新を用い、2つの補正は次のようにコンパクトに書ける：

\begin{align*}
\begin{aligned}
x_{t\to t-1}
&=\alpha_{t-1}\underbrace{\bigl(\hat{x}_0(x_t)+\eta_t^{\mathrm{data}}G_t^{\mathrm{data}}(x_t,c)\bigr)}_{\text{A. データ空間}} \\
&\quad-\sigma_{t-1}\sigma_t\underbrace{\bigl(\hat{s}(x_t)+\eta_t^{\mathrm{latent}}G_t(x_t,c)\bigr)}_{\text{B. ノイズ空間}}, \quad (8.4.2)
\end{aligned}
\end{align*}


ここで $\eta_t^{\mathrm{data}},\eta_t^{\mathrm{latent}}\ge 0$ はそれぞれデータ空間・ノイズ空間補正の強度を制御する。

\subsubsection*{A. データ空間でのガイダンス}


データ空間補正は、タスク損失を減少させる局所ステップを取ることで現在のクリーン推定を直接改善する：

\[
\hat{x}_0(x_t)\longmapsto\hat{x}_0(x_t)+\eta_t^{\mathrm{data}}G_t^{\mathrm{data}}(x_t,c).
\]


こうしてクリーン予測はまず条件とのより良い整合へ動かされ、次いで後続のサンプリング更新に用いられる。複数の最適化ステップが望ましいとき、この補正を繰り返し適用することもできる。
代表例には MGPD（He et al., 2023）と UGD（Bansal et al., 2023）が含まれる。

\subsubsection*{B. ノイズ空間でのガイダンス}


密接に関連する戦略は、サンプリング中に用いるスコアへ直接ガイダンスを適用する。§8.1で導入したポテンシャル視点に従い、共通の構成は現在のデノイズ推定上でタスク損失を評価し

\[
h_t(x_t;c):=\exp\bigl(-\tau\,\ell\bigl(\hat{x}_0(x_t),c\bigr)\bigr),\qquad\tau>0
\]


を選ぶ。上の $G_t$ の定義により、これは

\[
G_t(x_t,c)=\nabla_{x_t}\log h_t(x_t;c)=-\tau\nabla_{x_t}\ell\bigl(\hat{x}_0(x_t),c\bigr)
\]


を与える。したがって式 (8.4.2) のスコア項は

\[
\hat{s}(x_t)+\eta_t^{\mathrm{latent}}G_t(x_t,c)
\]


となる。この実務的更新を式 (8.1.3) の母集団水準構成と比較すると、$\eta_t^{\mathrm{latent}}$ が一般ガイダンス強度 $w_t$ の役割を果たし、$\hat{s}(x_t)$ が厳密な無条件スコアをその学習推定で置き換える。
$h_t(x_t;c)$ が真の条件付き尤度 $p_t(c|x_t)$ に、$x_t$ に独立な因子を除いて比例するとき、$G_t(x_t,c)$ はBayesガイダンス方向 $\nabla_{x_t}\log p_t(c|x_t)$ を回復する。一般のタスク損失に対しては、$G_t$ は代わりに代理ガイダンス方向を定義する。
$G_t$ の評価は $x$-予測を通じた微分を要する：

\[
G_t(x_t,c)=-\tau\bigl[\nabla_{x_t}\hat{x}_0(x_t)\bigr]^\top\nabla_{x_0}\ell(x_0,c)\bigm|_{x_0=\hat{x}_0(x_t)},
\]


これは実務上かなりの計算コストを生じ得る。
代表例には (Yu et al., 2023; Chung et al., 2022; Bansal et al., 2023) が含まれる。
上の構成は、ガイダンスが各サンプリングステップをどう修正するかを指定する。それ自体では、これらの修正されたダイナミクスがすべてのノイズ水準にわたって積分された後に得られる周辺分布を定めない。次にこの区別を明確にする。

\subsection*{8.4.2 誘導サンプリングはどの分布を生成するか？}


上のポテンシャル視点は、母集団水準で厳密な固定時刻密度解釈をもつ。特に、学習済み無条件スコア推定 $\hat{s}$ を厳密スコアで置き換えると、正のポテンシャル $h_t(x;c)$ とガイダンス強度 $w_t\ge 0$ が理想的な誘導場

\[
\nabla_x\log p_t(x)+w_t\nabla_x\log h_t(x;c)
\]


を定義する。各固定 $t$ で、これは

\[
\tilde{p}_t^{\mathrm{tilt}}(x|c)=\frac{1}{Z_t(c)}\,p_t(x)\,h_t(x;c)^{w_t},\qquad
Z_t(c)=\int p_t(x)\,h_t(x;c)^{w_t}\,dx \quad (8.4.3)
\]


のスコアである（$0<Z_t(c)<\infty$ を仮定）。実務で用いる学習場は、学習済みスコア推定 $\hat{s}(x_t)$ を通じてこの理想構成を近似する。
上の損失に基づく構成では

\[
h_t(x_t;c)=\exp\bigl[-\tau\,\ell\bigl(\hat{x}_0(x_t),c\bigr)\bigr]
\]


であり、$\eta_t^{\mathrm{latent}}$ が $w_t$ の役割を果たす。
この固定時刻恒等式は、誘導ダイナミクスがノイズ水準にわたって積分された後に生成される分布を定めない。したがって、一般の損失とガイダンススケジュールに対し、$T$ から $0$ へ誘導PF-ODEを解いても、一般には $\tilde{p}_0^{\mathrm{tilt}}(\cdot|c)$ に従って分布するサンプルは生じない。
理由はクロスタイム整合である：異なるノイズ水準での密度は独立に選べない。$\{\tilde{p}_t^{\mathrm{tilt}}\}_t$ が同じ前方拡散に付随する周辺経路をなすには、対応するFokker–Planck方程式も満たさねばならない：

\[
\partial_t\tilde{p}_t^{\mathrm{tilt}}+\nabla_x\cdot\bigl(f(t)x\,\tilde{p}_t^{\mathrm{tilt}}\bigr)-\frac{1}{2}g^2(t)\Delta_x\tilde{p}_t^{\mathrm{tilt}}=0.
\]


この条件が成り立つとき、対応するPF-ODE

\[
\frac{dx(t)}{dt}=f(t)x(t)-\frac{1}{2}g^2(t)\nabla_x\log\tilde{p}_t^{\mathrm{tilt}}\bigl(x(t)|c\bigr)
\]


は、対応する終端分布 $\tilde{p}_T^{\mathrm{tilt}}(\cdot|c)$ から初期化すれば $\tilde{p}_t^{\mathrm{tilt}}$ をその周辺経路としてもつ。
整合条件は直接確認できる。式 (8.4.3) をFokker–Planck方程式に代入し、$p_t$ 自体が同じ方程式を満たす事実を用いると

\begin{align*}
\begin{aligned}
&\partial_t\tilde{p}_t^{\mathrm{tilt}}+\nabla_x\cdot\bigl(f(t)x\,\tilde{p}_t^{\mathrm{tilt}}\bigr)-\frac{1}{2}g^2(t)\Delta_x\tilde{p}_t^{\mathrm{tilt}} \\
&=\tilde{p}_t^{\mathrm{tilt}}\Biggl[
\partial_t\log(h_t^{w_t})+f(t)x\cdot\nabla_x\log(h_t^{w_t})-g^2(t)\nabla_x\log p_t\cdot\nabla_x\log(h_t^{w_t}) \\
&\qquad-\frac{1}{2}g^2(t)\Bigl(\Delta_x\log(h_t^{w_t})+\bigl\|\nabla_x\log(h_t^{w_t})\bigr\|^2\Bigr)-\frac{Z_t'}{Z_t}
\Biggr]
\end{aligned}
\end{align*}


が得られ、引数 $(x;c)$ は読みやすさのため省略する。したがって傾いた族が妥当な周辺経路であるのは、括弧内の量が消えるときに限る。一般のタスク損失とガイダンススケジュールはこの条件を満たす必要がない。したがって不一致は、厳密なスコア評価と厳密なODE積分の下でも残り得る：それは数値離散化ではなく誘導ダイナミクス自体に由来する。これは (Lai et al., 2023a) で議論される、より広いFokker–Planck整合要件の一例である。
2つの単純な場合が区別を明確にする。$w_t\equiv 0$ のとき傾きは消え、元の無条件経路 $p_t$ が回復される。厳密なクラシファイア・ガイダンスでは、前方拡散が $x$ を壊しつつ $c$ を固定する標準設定の下で

\[
h_t(x;c)=p_t(c|x),\qquad w_t\equiv 1.
\]


Bayesの公式により

\[
\tilde{p}_t^{\mathrm{tilt}}(x|c)=p_t(x|c)
\]


となり、傾いた族はちょうど条件付き周辺経路である。これがオラクルクラシファイア尤度に対するガイダンス重み $1$ の特別な役割を説明する。
一般のガイダンス強度、例えば

\[
p_t(x)\,p_t(c|x)^\omega
\]


では、表現は正規化可能な限りなお妥当な固定時刻傾き密度を定義するが、一般に誘導サンプラの周辺を記述することは保証されない。特別な選択は必要な整合条件を満たし得るので、$\omega\neq 1$ であることだけでは不整合を含意しない。
Feynman–Kac Correctors（skreta2025feynman）のような補正に基づくアプローチは、規定された傾き分布列をより忠実に追うことが目標のとき、このクロスタイム不一致を明示的に扱う。

\subsection*{8.4.3 逆問題への訓練不要アプローチの例}


§8.4.1で導入した原理は逆問題に重要な応用をもつ。まず背景の概観から始め、次いで事前学習済み拡散モデルを推論時の逆問題求解にどう活用するかを示すいくつかの具体例を与える。

\textbf{逆問題の背景。} $A$ を既知の前方演算子（非線形でもよい）とし、

\[
y=A(x_0)+\sigma_y z,\qquad z\sim\mathcal{N}(0,I) \quad (8.4.4)
\]


とする。Bayes目的は事後 $p_0(x_0|y)$ を特徴づけることであり、これは測定モデルと妥当なクリーン信号に関する事前知識を組み合わせる。ブラインド逆問題はこの設定を拡張し、$A$ の未知パラメータを追加の推論量として扱う。
古典的逆問題手法は、変分またはBayes定式化において、データ忠実項と解析的または手設計の事前を組み合わせることが多い。教師あり復元手法は代わりに、ペアデータで測定からクリーン信号への直接写像を学習し得る。事前学習済み拡散モデルは第3の経路を与える：クリーンデータ上の学習済み事前を測定演算子をまたいで再利用し、測定情報は推論時にのみ導入する。

\textbf{逆問題ソルバとしての事前学習済み拡散モデル。} 上の測定モデルに対し、Bayesの公式は

\[
\nabla_{x_t}\log p_t(x_t|y)=\underbrace{\nabla_{x_t}\log p_t(x_t)}_{\text{データスコア}}+\underbrace{\nabla_{x_t}\log p_t(y|x_t)}_{\text{測定ガイダンス}} \quad (8.4.5)
\]


を与える。データスコアは事前学習済み拡散モデルにより供給できる一方、測定項は一般に閉形式では得られない。したがって訓練不要の逆問題ソルバはサンプリング中に測定情報を導入するが、すべてが同じ数学的形でそうする必要はない。
有用な区別は2つの共通戦略の間にある。いくつかの手法は中間サンプルを修正または射影することで測定整合を直接強制する。他は尤度スコア $\nabla_{x_t}\log p_t(y|x_t)$ を近似し、式 (8.4.5) を通じてデータスコアに加える。

\textbf{ILVR（Choi et al., 2021）。} ILVRは第1戦略の単純な例を与える。$\varphi_N$ をその線形低域通過演算子とする。各逆ステップで、まず無条件提案

\[
x_{t-1}'\sim p_\theta(x_{t-1}'|x_t)
\]


をサンプリングし、参照画像を同じ水準へノイズ化する：

\[
y_{t-1}\sim p_{t-1}(y_{t-1}|y).
\]


提案は次いで

\[
x_{t-1}=\varphi_N(y_{t-1})+(I-\varphi_N)x_{t-1}'
\]


に従って洗練される。こうして規定された低周波成分はノイズ化された参照から来て、相補成分は無条件拡散提案から来る。これは $\nabla_{x_t}\log p_t(y|x_t)$ の明示的近似ではなく、サンプル空間の整合補正である。

\textbf{Diffusion Posterior Sampling（DPS）（Chung et al., 2022）。} 加法ガウス測定ノイズ $\sigma_y>0$ をもつ既知の非線形前方演算子 $A$ に対し、Diffusion Posterior Sampling（DPS）は

\[
\nabla_{x_t}\log p_t(y|x_t)\approx\nabla_{x_t}\log p\bigl(y\bigm|X_0=\hat{x}_0(x_t)\bigr) \quad (8.4.6)
\]


を近似する。ここで $\hat{x}_0(x_t):=\mathbb{E}[x_0|x_t]$ は $x_t$ が与えられたときのクリーンサンプルの事後平均であり、典型的にはTweedieの公式（式 (3.3.6)）を用いて事前学習済み拡散モデルから推定される。
このプラグイン近似は、測定尤度の事後平均を事後平均での値で置き換える：

\begin{align*}
\begin{aligned}
p_t(y|x_t)
&=\int p_t(y|x_t,x_0)\,p(x_0|x_t)\,dx_0 \\
&=\int p_t(y|x_0)\,p(x_0|x_t)\,dx_0
\approx p_t\bigl(y\bigm|X_0=\hat{x}_0(x_t)\bigr),
\end{aligned}
\end{align*}


ここで $y$ は $x_0$ が与えられれば（$x_t$ にではなく）$x_0$ のみに依存することを用い、近似は事後 $p(x_0|x_t)$ がその平均の周りに強く尖っているという仮定の下で成り立つ。
\[
p_t\bigl(y\bigm|X_0=\hat{x}_0(x_t)\bigr)=\mathcal{N}\bigl(y;A(\hat{x}_0(x_t)),\sigma_y^2 I\bigr)
\]

なので、

\begin{align*}
\begin{aligned}
\nabla_{x_t}\log p_t(y|x_t)
&\approx\nabla_{x_t}\log\mathcal{N}\bigl(y;A(\hat{x}_0),\sigma_y^2 I\bigr) \\
&=-\frac{1}{2\sigma_y^2}\nabla_{x_t}\|y-A(\hat{x}_0)\|^2 \\
&=\frac{1}{\sigma_y^2}\Bigl[J_A\bigl(\hat{x}_0(x_t)\bigr)\cdot\nabla_{x_t}\hat{x}_0(x_t)\Bigr]^\top\bigl(y-A(\hat{x}_0(x_t))\bigr)
\end{aligned}
\end{align*}


を計算する。ここで $J_A\bigl(\hat{x}_0(x_t)\bigr):=\nabla_{x_0}A(x)\big|_{x=\hat{x}_0(x_t)}$ は前方演算子の入力に関するヤコビアンを表す。この公式はスコア近似パイプラインを通じて勾配を伝播し、ノイズ付きサンプル $x_t$ の摂動に対して測定尤度がどう変化するかを反映する。
線形逆問題では、これはさらに

\[
\nabla_{x_t}\log p_t(y|x_t)\approx\frac{1}{\sigma_y^2}\bigl[A\cdot\nabla_{x_t}\hat{x}_0(x_t)\bigr]^\top\bigl(y-A(\hat{x}_0(x_t))\bigr)
\]


に簡約される。
$\nabla_{x_t}\log p_t(y|x_t)$ の様々な近似を提案する拡散に基づく逆問題ソルバの研究は多数ある。包括的概観は Daras et al. (2024) のサーベイを参照されたい。

\subsubsection*{関連する教師あり反復復元視点とPF-ODEへの接続}


上の議論は、事前学習済み拡散事前と式 (8.4.5) のBayes分解を組み合わせて逆問題を解く。より具体的には、これらの手法は事前分布のスコアと測定モデルを用いて再構成過程を誘導する。同時に、より広い着想も反映する：クリーン信号を1ステップで回復するのではなく、多くの小さな更新を通じて壊れた入力を徐々に改善する。類似の高水準戦略は教師あり画像復元にも現れるが、出発点は異なる。事前学習済みスコアモデルを事前として用いるのではなく、これらの手法はペアの壊れた例とクリーン例から洗練手続きを直接学習する。
代表例は InDI（delbracio2023inversion）である。上のBayes則に基づくアプローチとは異なり、InDIはペアデータ $(y,x_0)$——$y$ が壊れた観測、$x_0$ が対応するクリーンターゲット——から始まる。重要なことに、一般には $x_0$ から $y$ への破損が解析的に既知であることやガウスであることを仮定しない。
基本的着想は、2つの端点を補間経路

\[
x_t=(1-t)x_0+t y,\qquad t\in[0,1]
\]


で結ぶことである。この経路は、クリーンターゲットから壊れた観測への便利な軌道としてのみ理解すべきである。一般にはガウスノイズ化過程ではない。この経路に沿って、任意の $0\le s\le t$ に対し厳密な関係

\[
x_s=\Bigl(1-\frac{s}{t}\Bigr)x_0+\frac{s}{t}x_t
\]


が成り立つ。$x_t$ が与えられた条件付き期待値をとると

\[
\mathbb{E}[x_s|x_t]=\Bigl(1-\frac{s}{t}\Bigr)\mathbb{E}[x_0|x_t]+\frac{s}{t}x_t
\]


が得られる。これは、現在状態 $x_t$ からわずかに劣化の少ない状態へ動くには、$x_t$ からクリーンターゲットを予測すれば十分であることを示す。
この着想に基づき、ペア例から時間条件付き予測器 $F_\phi(x_t,t)$ を、回帰目的

\[
\min_\phi\,\mathbb{E}_{(x_0,y)\sim p(x_0,y),\,t\sim p(t)}\bigl\|F_\phi\bigl((1-t)x_0+ty,\,t\bigr)-x_0\bigr\|_2^2
\]


を最小化して訓練する。ここで $p(x_0,y)$ はクリーンターゲット $x_0$ とその壊れた観測 $y$ からなるデータ対の結合分布、$p(t)$ は訓練中に補間時間をどうサンプリングするかを指定する。すなわち $F_\phi$ は、補間経路上の中間劣化した状態からクリーンターゲットを予測するよう訓練される。この二乗誤差損失に対し、Bayes最適予測器は条件付き期待値

\[
F^*(x_t,t)=\mathbb{E}[x_0|x_t]
\]


である。したがってモデルクラスが十分表現力豊かで最適化が成功すれば、学習済み予測器は時刻 $t$ の中間状態 $x_t$ に付随する平均クリーンターゲットを近似する。この意味で、学習目的は拡散モデル訓練と精神的に類似する：いずれの場合も、中間の壊れた状態からよりクリーンなターゲットへの時間条件付き回帰を学習する。違いは、拡散モデルが規定された前方破損過程から合成的にそのような訓練対を生成するのに対し、ここでは端点対 $(y,x_0)$ が教師あり復元データから直接来ることである。
推論時には $\hat{x}_1=y$ から出発し、

\[
\hat{x}_{t-\Delta t}=\frac{\Delta t}{t}F_\phi(\hat{x}_t,t)+\Bigl(1-\frac{\Delta t}{t}\Bigr)\hat{x}_t
\]


により推定を繰り返し洗練する。ゆえに最終再構成は単一の直接予測ではなく段階的洗練により生じる。連続時間極限では、これは InDI の残差フローODE

\[
\frac{dx_t}{dt}=\frac{x_t-F_\phi(x_t,t)}{t}
\]


へ至る。
さらに、破損が既知の解析形

\[
y=x_0+\sigma\epsilon,\qquad\epsilon\sim\mathcal{N}(0,I)
\]


（既知ノイズ水準 $\sigma$）をとると仮定するガウスデノイジング設定へ特殊化すると、InDIの残差フローODEは厳密に確率流ODEへ還元される。実際、補間経路は

\[
x_t=(1-t)x_0+ty=x_0+t\sigma\epsilon
\]


となり、これはまさに標準的な拡散形

\[
x_t=\alpha_t x_0+\sigma_t\epsilon,\qquad\alpha_t=1,\ \sigma_t=t\sigma
\]


である。したがって予測器 $F_\phi(x_t,t)\approx\mathbb{E}[x_0|x_t]$ は、ガウス破損した中間状態 $x_t$ をデノイズするよう学習している。この場合、残差フローODEは

\[
\frac{dx_t}{dt}=\frac{x_t-\mathbb{E}[x_0|x_t]}{t}=-\dot{\sigma}_t\sigma_t\nabla_{x_t}\log p_t(x_t)
\]


となり、最後の等号は $\sigma_t=t\sigma$ とTweedieの公式

\[
\mathbb{E}[x_0|x_t]=x_t+\sigma_t^2\nabla_{x_t}\log p_t(x_t)
\]


から従う。これはこのガウスデノイジング経路に対するPF-ODEそのものであり、デノイザ形またはスコア形で書かれる。こうしてガウスデノイジング特殊化の下で、InDIは拡散モデリングの別解釈を示唆する：ガウス破損データ上の教師ありデノイザ学習とみなし、PF-ODEは段階的洗練の連続時間極限として生じる、という見方である。
\end{document}
